\needspace{5cm} \section{Abeceda} \label{abeceda} \katabox{Téma 
}{2 
}{2 min} 
 
V~kategorii \textbf{Abeceda} vám 2 hráči zahrají na Vaše téma, ale tak, že jejich repliky musí začínat písmeny tak, jak jsou seřazeny v~abecedě. 
 
\subsection{Průběh} Připraví se dva hráči a je jim zadáno téma či prostředí. Hráči poté rozehrají improvizaci, ale s~omezením, že jejich repliky musí začínat písmeny podle pořadí v~abecedě. V~případě, že se hráč splete může být nahrazen jiným hráčem ze svého týmu. Pokud hráči dojdou až na konec abecedy, začíná se opět od písmene A. 
 
\subsection{ Varianty } \begin{itemize}
\item  Hraje se i na písmena s~háčky a čárkami.
\item  Mohou se zapojit další hráči.
\item  Kategorie končí písmenem Z.
\end{itemize}
 
\subsection{ Běžné chyby } \begin{itemize}
\item Hráč se ztratí v~abecedě.
\end{itemize}
 
 
 
 
 
\needspace{5cm} \section{Báseň (cvičení)} \label{báseň (cvičení)} Na zadané téma diváků, které může být buď inspirací nebo přímo názvem \odkaz{básně}{báseň}, hráči postupně chodí dopředu a každý přednáší  
jeden verš. Verše se rýmují \odkaz{ABAB}{abab}. 
 
V~případě více slok, hráči, kteří složili 1. sloku opouští hrací prostor a vrací se k~ostatním hráčům a kdokoliv z~nich může začít tvořit 
další sloku. 
 
\subsection{Varianty} \begin{itemize}
\item  báseň končí po po první sloce.
\end{itemize}
  
\subsection{Viz také} \begin{itemize}
\item \odkaz{Poetická}{poetická}
\item \odkaz{Divadlo poezie}{divadlo poezie}
\item \odkaz{Báseň od publika}{báseň od publika}
\end{itemize}
 
\subsection{Běžné chyby} \begin{itemize}
\item Důležitým prvkem u~básně je zvýraznit poslední řádek, dát mu důraz, aby posluchač cítil, že touto větou báseň končí.
\item Tím, že verše vznikají na jevišti a nejistotou hráčů, mohou básně působit občas plitce. Je tedy třeba dbát v~tomto cvičení na přednes a sebevědomí hráčů.
\end{itemize}
 
 
 
 
\needspace{5cm} \section{Break-up song} \label{break-up song} \katabox{Dobrovolník z~publika, židle 
}{2 
}{neomezený, 3 - 5 minut (podle počtu hráčů)} 
 
Hráči zazpívají o~rozchodu na motivy informací od jednoho dobrovolníka. 
 
\subsection{ Průběh } Průběh je téměř stejný jako u~kategorie \odkaz{Love song}{love song}. \odkaz{MC}{mc} ale z~publika požádá diváka, který v~nedávné době prošel rozchodem. Píseň je mířená na jeho protějšek. 
 
\subsection{ Varianty } \begin{itemize}
\item \odkaz{Love song}{love song}
\end{itemize}
 
 
 
 
 
\needspace{5cm} \section{Divadlo poezie} \label{divadlo poezie} \katabox{Téma 
}{4 a více 
}{4 minuty} 
 
 
\textbf{Divadlo poezie} poetickým způsobem ztvární vaše téma. 
 
 
\subsection{Průběh} Hráči společně tvoří pomocí zvuků, rýmu a pohybu zvukomalebný příběh. 
Využívají přitom opakování slov, zvuků, pohybů, celá skupina se organicky přelévá a využívá synonyma, asociace a přirovnání. 
 
{{todo|rozšířit text, přidat videa.}} 
 
 
 
 
 
\needspace{5cm} \section{Dopředu, dozadu} \label{dopředu, dozadu} \katabox{Téma 
}{2 a naskakující 
}{neomezený} 
 
V~této kategorii se podíváme na příběh několikrát a zajímavé části si přetočíme dopředu nebo zpátky. 
 
\subsection{ Průběh } Hraje se volně příběh. \odkaz{MC}{mc} může kdykoliv přepnout směr děje příběhu dopředu nebo dozadu. Pokud \odkaz{hráči}{hráč} hrají příběh směrem dopředu a zazní povel \uv{Dozadu!} , opakují pozpátku všechny pohyby a repliky v~opačném pořadí než zazněly. Řeč se při směru dozadu neobrací, mluví se normálně. Příběh se odehrává až na začátek nebo dokud \odkaz{MC}{mc} opět nepřepne směr. Poté, co \odkaz{MC}{mc} přepne opět dopředu, všechno se zopakuje, dokud se \odkaz{hráči}{hráč} nedostanou k~bodu, kde ještě příběh nebyl zahrán a pokračují v~ději. 
 
\subsection{ Běžné chyby } \begin{itemize}
\item  Příliš mnoho povídání, statická scénka
\item  Hráči zapomenou, co se v~příběhu stalo
\item  \odkaz{MC}{mc} pouští příliš dlouhé úseky dozadu
\end{itemize}
 
 
 
 
 
\needspace{5cm} \section{Dva začátky} \label{dva začátky} \katabox{Téma 
}{libovolný 
}{5 minut} 
 
 
Diváci si ze \textbf{dvou začátků} příběhu vyberou, který chtějí slyšet. 
 
 
\subsection{Průběh} Dva hráči krátce shrnou příběh, přednesou anotaci či ukázku - dva rozličné příběhy na stejné téma. 
Následně se \odkaz{hlasováním}{hlasování} určí, který příběh bude pokračovat, ten následně hráči sehrají. 
 
\subsection{ Varianty } \begin{itemize}
\item  \odkaz{Režiséři}{režiséři}
\end{itemize}
 
 
 
\needspace{5cm} \section{Dvě repliky} \label{dvě repliky} \katabox{Dvě věty pro kažého (až na jednoho) hráče 
}{libovolný (3 a více) 
}{3 minuty} 
 
Hráči mají jen \textbf{dvě repliky} na to, aby vyjádřili vše, co chtějí říci. 
 
\subsection{Průběh} Jeden hráč není ve svém vyjadřování nijak omezován, další hráči mají přidělen každý dvě jiné obecné věty (např. Nesnáším tě.  Tohle by se za Marie Terezie nestalo.). Jenom \odkaz{replikami}{replika} mohou reagovat na jakékoli promluvy centrálního hráče, který své nabídky volí tak, aby ostatním tyto reakce umožnil. 
 
\subsection{Běžné chyby} \begin{itemize}
\item  Neomezovaný hráč nenahrává těm ostatním.
\item  Příliš obskurní věty pro hráče.
\end{itemize}
 
 
 
\needspace{5cm} \section{Facebookový status} \label{facebookový status} \katabox{Poslední facebookový status některého diváka 
}{libovolný  
}{3  minuty} 
 
Hráči sehrají příběh ze života inspirovaný  \textbf{Facebookovým statusem} některého diváka. 
 
\subsection{Průběh} Hraje se příběh statusem inspirovaný, situace která končí tímto statusem.  
 
 
 
\needspace{5cm} \section{Handicapovaná pohádka} \label{handicapovaná pohádka} \katabox{Známá pohádka a handicap pro každého hráče 
}{5 
}{podle vybrané pohádky} 
 
 
Modifikace známé pohádky od diváků pomocí handicapů. 
 
\subsection{Průběh}  
Hraje obvykle 5 hráčů, každý z~nich má jeden handicap (verbální, fyzický, mentální,...) a společně hrají známou pohádku zadanou od diváků. Rozhodčí předem určuje handicapy a jednotlivé postavy, které hráči hrají. 
Pohádku musí všichni hráči znát, v~opačném případě je vhodné např. nechat diváka odříkat celou pohádku. 
 
\subsection{Varianty} \begin{itemize}
\item Pohádka s~charakterovou vadou - hráči mají zadanou pro svou postavu výraznou emoci/charakterovou vlastnost (podobně jako do \odkaz{Schizofrenie}{schizofrenie})
\end{itemize}
 
 
 
 
\needspace{5cm} \section{Jestli víš, co tím myslím} \label{jestli víš, co tím myslím} \katabox{Prostředí 
}{2+  
}{2  minuty} 
 
Každé slovo může získat zcela nový význam, \textbf{jestli víte, co tím myslím}. 
 
\subsection{Průběh} Dva hráči rozehrávají scénu v~daném prostředí, ale každou svou \odkaz{repliku}{replika} zakončí  
větou jestli víš, co tím myslím. Všechny věty tak získávají jakousi \uv{dvousmyslnost} , která se ještě podpoří volbou slov před tímto zakončovačem repliky.   
 
 
 
\needspace{5cm} \section{Kaskadéři} \label{kaskadéři} \katabox{Téma 
}{4 (dva páry) 
}{orientačně 3 min} 
 
Přepínání mezi akčním a všedním životem 
 
\subsection{ Průběh }  
Hrají 4 hráči, najednou ale vždy jen dva. Jedna dvojice hraje civilní příběh na téma diváků, ale jakmile rozhodčí zakřičí “kaskadéři”, tak se vymění za druhou dvojici, která pokračuje ve stejném příběhu a stejných postavách, ale vše je akční, přeháněné a zběsilé, jakoby kaskadéři zaskočili za ty náročné části příběhu. Po chvíli rozhodčí zakřičí “civil” a vrací se první dvojice. Takto se civilní dvojice a kaskadérská několikrát během improvizace vystřídají. 
 
\subsection{ Běžné chyby } \begin{itemize}
\item  Ztráta napojení s~druhou dvojicí
\end{itemize}
 
 
 
\needspace{5cm} \section{Kramářská píseň} \label{kramářská píseň} \katabox{Téma 
}{libovolný (4 a více) 
}{4 minuty} 
 
\textbf{Kramářská píseň} sloužila předkům jako bulvární zdroj informací. 
 
\subsection{Průběh} Jeden hráč se stává zpěvákem, na jeho tlesknutí se velmi rychle mění \odkaz{živý obraz}{živé obrazy} tvořený ostatními hráči. 
Hráč zpěvem popisuje děj, a obrazy, mezi slokami obvykle přechází do vyvolavačské komunikace s~diváky ( Ale to ještě není vše, přijďte blíž a uvidíte, přijďte blíž a uslyšíte.) 
 
Pro zjednodušení této obtížné kategorie se občas používá melodie z~\uv{Pes jitrničku sežral}  
 
\subsection{Varianty} \begin{itemize}
\item  Na zápase mohou být dva zpěváci, kteří se střídají po sloce a promluvě.
\end{itemize}
 
\subsection{Viz také} Podobná je kategorie \odkaz{Diapozitivy}{diapozitivy}. 
 
 
 
 
 
\needspace{5cm} \section{Muzikál} \label{muzikál} \katabox{Slovní téma 
}{2 - 5 
}{5-15 minut} 
 
Muzikál je improvizační kategorie založena na příběhu (pro lepší přehlednost zpravidla jednoho hlavního hrdiny) do kterého jsou vkládány písně, které posunují děj vpřed. Důležitá a hojně opomíjena je složka taneční, choreografická. Není ovšem nutné stále jen zpívat (narozdíl od kategorie \odkaz{opera}{opera}) v~muzikálu jsou stejně nutné i dialogy. 
 
\subsection{Jak na to}  
Důležité je mít příběh (nejlépe hlavní hrdina někam musí jít a něco udělat aby něco se stalo). Často se hledí hlavně na zpěv, ten však je pouze prostředek ke sdělení nějaké informace (proto když vám prostě nejde rým tak tam nerýmujte, ale veďte příběh dál). Tip na tvoření choreografií: postavě se tak, aby jeden byl kousek před ostatními ten pak udává pohyby které ti za ním zrcadlí (vytváří tak efekt nacvičených choreografií) 
 
{{Todo| 
\begin{itemize}
\item  Doplnit průběh
\item  Doplnit video
\end{itemize}
}} 
 
 
 
 
 
\needspace{5cm} \section{Občanky} \label{občanky} \katabox{Občanské průkazy od diváků podle počtu hráčů 
}{3-4 
}{volný, dokud nedojdou nápady} 
 
Kategorie, ve které herci ztvární postavy vybraných diváků. 
 
\subsection{Průběh}  
Hraje 4-5 hráčů, každý si z~diváků vybere jednoho člověka, který má stejné pohlaví, jako hráč. Divák mu půjčí občanku, řidičák, nebo pas a hráč má pak krátký čas na “nastudování” informací, jelikož poté hraje v~improvizaci právě toho diváka, kterého si vybral. Je zajímavépřejmout i charakteristiky/vystupování postavy, kterou hrajeme. 
 
 
 
\needspace{5cm} \section{Oprsklá} \label{oprsklá} \katabox{téma/prostředí/vztah/situace 
}{2 
}{neomezen 
} 
 
V~kategorie \textbf{Oprsklá} jsou hráči opravdu překvapeni každou informací, kterou se dozví. 
 
\subsection{Příprava} Každý hráč obdrží naplněnou sklenici vody.  
 
\subsection{Průběh} Kategorie spočívá v~tom, že si navzájem sdělují informace, které druhého překvapí. Výrazem překvapení je následné vyprsknutí vody na spoluhráče. Hraje se, dokud je v~půllitru nějaká voda.  
 
\subsection{Úklid} Po kategorii je scéna i hráči mokrá. Doporučujeme uklidit. 
 
\subsection{Běžné chyby} \begin{itemize}
\item Hráč ve své větě, větách nepřinese žádnou informaci.
\item Hráč nereaguje udiveně.
\item Hráč neumí vodu vyprsknout (doporučujeme natrénovat)
\end{itemize}
 
\subsection{Viz také} Kategorie je částečně podobná \odkaz{Co tím chceš jako říct}{co tím chceš jako říct}, hodí se zejména na horké letní festivaly. 
 
 
 
\needspace{5cm} \section{Orákulum} \label{orákulum} \katabox{otázka 
}{všichni 
}{volný} 
 
V~této kategorii se z~hráčů stává mýtické všeznalé stvoření, kterého se diváci ptají na záludné otázky. 
 
\subsection{ Průběh } Hráči se k~sobě přitisknou a dohromady odpovídají na otázky z~publika. Diváci se ptají na cokoliv, často i řečnické otázky. Orákulum by mělo být schopné na ně originálně a jednohlasně odpovědět. 
 
\subsection{ Běžné chyby } \begin{itemize}
\item Dlouhé odmlky
\item Všichni mluví přes sebe
\item Orákulum má viditelně jednoho hlavního mluvčího, po kterém všichni opakují
\end{itemize}
 
 
 
 
\needspace{5cm} \section{Příběh po slovech} \label{příběh po slovech} \katabox{Téma 
}{libovolný,minimálně 2 
}{volný} 
 
 
Napojení improvizátoři vypráví příběh po slovech. 
 
 
\subsection{ Průběh }  
Herci stojí vedle sebe v~řadě a po jednom slově vyprávějí příběh. 
 
Jedno ze základních cvičení, které se používá jako součást kategorií \odkaz{Dopis}{dopis} nebo \odkaz{Divadlo poezie}{divadlo poezie}, ale také samostatně. 
 
 
\subsection{ Varianty }  
\begin{itemize}
\item  Diváci zadají ponaučení místo tématu. Příběh skončí přesnou citací zadaného ponaučení.
\end{itemize}
 
\subsection{ Běžné chyby }  
\begin{itemize}
\item  Přemýšlení nad dalším slovem a sebehodnocení
\item  Přehnaná kreativita
\end{itemize}
 
 
 
 
 
 
 
\needspace{5cm} \section{Režiséři} \label{režiséři} \katabox{Téma pro každého z~režisérů}{neomezen}{neomezen} 
 
V~kategorii \textbf{Režiséři} též známé jako \textbf{Super scene} několik hráčů v~roli režiséra představí divákům své filmy. Diváci mají poté možnost rozhodovat o~tom, který film chtějí vidět. vyberte si, který chcete vidět. 
 
\subsection{Průběh} Vyčlení se hráči, kteří se stávají režiséry. Každému z~režisérů je zadán název jeho filmu (téma od diváků) a jeho žánr. Každý režisér má poté max. 30s na krátkou slovní upoutávku tohoto filmu, kterým se snaží diváky navnadit. Po představení všech filmů diváci hlasují o~tom, který film vidět vůbec nechtějí. Z~ostatních filmů se poté odehraje jejich první část, během níž může režisér komentovat děj filmu zákulisními informacemi nebo děj přímo ovlivňovat. Každý režisér první část svého filmu ukončí. Ke slovu se opět dostává \odkaz{MC}{mc} a diváci. Opět je zde hlasování po kterém opět jeden z~filmů vypadne. Takto kategorie pokračuje až do konce jediného celého filmu. 
 
 
\subsection{ Varianty } \begin{itemize}
\item  \odkaz{Dva začátky}{dva začátky}
\end{itemize}
 
 
 
\needspace{5cm} \section{Scénář / Souboj s~knihou} \label{scénář / souboj s~knihou} \katabox{Připravená kniha či scénář 
}{2 
}{3 minuty} 
 
 
V~této  kategorii  vám hráči sehrají příběh, který bude částečně shodný s~některým známým dílem. 
 
\subsection{Varianty začátku} \begin{itemize}
\item \textbf{Scénář} - Jeden hráč dostane od \odkaz{MC}{mc} přidělený kus scénáře - dialogu, ve kterém je vše co říká jedna postava vyškrtáno. Hráč musí používat pouze věty z~dialogu té druhé postavy a to v~původním pořadí.
\item \textbf{Souboj s~knihou} - Jeden hráč dostane přidělenou knihu.  Všechny věty, které pronese, musí být \odkaz{repliky}{replika} z~této knihy.
\end{itemize}
 
\subsection{Průběh} Sehraje se scénka, kdy druhý hráč musí citlivě reagovat a pečlivě poslouchat prvního hráče, v~případě Scénáře navíc odhadovat, co se asi v~scénáři může stát, aby reakce druhého hráče vypadaly logicky. 
 
\subsection{Viz také} \begin{itemize}
\item  \odkaz{Dvě repliky}{dvě repliky}
\end{itemize}
 
 
 
\needspace{5cm} \section{Shakespeare pozpátku} \label{shakespeare pozpátku} \katabox{Téma 
}{neomezený 
}{cca 10 minut} 
 
Uvidíme pravé a hutné drama jak od Shakespeara, ale celé \textbf{pozpátku}. 
\subsection{Průběh} Hraje se příběh rozdělený do 5 dějství pojmenovatelných dle modifikovaných částí \odkaz{příběhu}{příběh}.  
Tyto dějství se ale hrají v~opačném pořadí - začínáme tedy od konce.  
\begin{itemize}
\item Tragedie
\item Konflikt/Intriky
\item Zápletka
\item Láska/Vztahy
\item Úvod
\end{itemize}
 
V~rámci dějství jsou scény řazeny standardně, chronologicky.  
 
 
 
\needspace{5cm} \section{Statusy} \label{statusy} \katabox{Téma a lístky s~čísly 
}{4 
}{volný} 
 
 
Hádání statusu. 
 
 
\subsection{ Průběh }  
Hrají 4 hráči a každý má na čele vylosovaný lístek s~číslem (1-10), který ale neviděl a vylosoval si ho poslepu. Poté se zadá prostředí, kde se často vyskytují lidé s~různými společenskými postaveními a přichází první hráč. Ten by měl začít neutrální činností, jelikož neví, jaké číslo (jaký status) má. Obecně platí, že čím vyšší číslo, tím vyšší společenské postavení. Hráči po skončení improvizace tipují, jaké číslo mají na čele. 
 
 
 
\needspace{5cm} \section{Stojí, leží, sedí} \label{stojí, leží, sedí} \katabox{Téma}{3}{2 min} 
 
V~Kategorii \textbf{stojí, leží, sedí} hrají 3 hráči a vždy musí jeden z~hráčů stát, jeden sedět a poslední ležet. 
 
\subsection{ Průběh } Jsou vybrání tři hráči, je jim zadáno téma a dána k~dispozici jedna židle. Hráči zaujmou dané pozice - jeden zůstane stát, druhý si sedne a třetí lehne. Poté rozehrají improvizaci na dané téma a přitom musí zachovávat, aby byly stále zastoupeny všechny tři pozice. Hráči mohou své pozice během improvizace měnit. Není určeno pořadí v~jakém musí pozice jít po sobě. Ležící hráč si může rovnou stoupnout a stejně tak i obráceně.  
 
\subsection{ Časté chyby } \begin{itemize}
\item Hráči nedávají pozor a opomenou jednu pozici
\end{itemize}
 
 
 
 
\needspace{5cm} \section{Tříhlavý song} \label{tříhlavý song} \katabox{první verš básně}{3}{cca 3 min} 
V~této takřka nehrané kategorii hráči zazpívají píseň po jednom slově. 
\subsection{Průběh} Hráči se k~sobě přitisknou a dohromady zazpívají píseň s~tím, že se střídají po jednom slově. První verš, vyžádaný od publika, na začátku zopakují, čímž se i dostanou do nějakého tempa.  
 
\subsection{Běžné chyby} \begin{itemize}
\item Dlouhé odmlky
\item Běžné veršovací chyby
\end{itemize}
 
\subsection{Viz také} \begin{itemize}
\item Kategorie \odkaz{Opilecká píseň}{opilecká píseň} - zpěv se strřídáním po jednom verši.
\item Cvičení \odkaz{Příběh po jednom slově}{příběh po jednom slově}
\item Rozcvička \odkaz{Vtip}{vtip}
\end{itemize}
 
 
 
 
\needspace{5cm} \section{Tři židle} \label{tři židle} \katabox{není nutné 
}{3 
}{5-15 minut} 
 
Každá skupinka stojících hráčů má svůj příběh. 
 
\subsection{Průběh} Připraví se 3 židle, na které se neutrálně usadí stejný počet herců. 
Hráči dle svého rozhodnutí vstávají, vždy však stojí alespoň jeden hráč. Každá možná podskupinka vypráví nějaký příběh, jeho vyprávění je však dočasně přerušeno, když se postaví nebo si sedne některý z~hráčů.   
Kategorie končí v~momentě, kdy jsou všechny příběhy dovyprávěny, hráči se proto snaží je dovyprávět cv podobný okamžik. 
 
\subsection{Princip} Kategorie je divácky zábavná díky změnám dynamice vyprávění. Herci nemají žádné \odkaz{askfor}{askfor}, mají proto úplnou svobodu ve výběru  žánru, formy a délce příběhů. 
 
\subsection{Varianty} \begin{itemize}
\item  Dvě židle (3 příběhy)
\item  Čtyři židle (15 příběhů)
\item   Pět židlí (31 příběhů)
\end{itemize}
 
\subsection{Podobné} \begin{itemize}
\item  \odkaz{Spoon river}{spoon river}
\item  \odkaz{Toaster}{toaster}
\end{itemize}
 
\subsection{ Může se plést s~} \begin{itemize}
\item  \odkaz{Židle}{židle}
\item  \odkaz{Dvě židle}{dvě židle}
\end{itemize}
 
 
 
 
 
 
 
\needspace{5cm} \section{Věty, které by neměly zaznít} \label{věty, které by neměly zaznít} \katabox{Situace/činnosti (3 - 5) 
}{libovolný 
}{neomezený} 
 
 
Hráči říkají věty, které by v~určitých situacích nebo na určitých místech neměly zaznít. 
 
 
\subsection{ Průběh } Podobně jako v~kategorii \odkaz{1000 způsobů jak}{1000 způsobů jak} vybere \odkaz{MC}{mc} od publika několik situací např. první kroky dítěte, přijímací pohovor, na operačním sále, atd. \odkaz{Hráči}{hráč} se postaví do půlkruhu a poté vždy ten s~nápadem na větu vyběhne na forbínu, kde tuto \odkaz{repliku}{replika} prezentuje a opět se zařadí do půlkruhu. Situace či prostředí se nechává do vyčerpání nápadů nebo na pokyn \odkaz{MC}{mc}. 
 
\subsection{ Varianty } \begin{itemize}
\item  \odkaz{1000 způsobů jak}{1000 způsobů jak}
\item  \odkaz{Metafory}{metafory}
\item  Některé věty se hodí do více situací a hráči toho hojně využívají.
\end{itemize}
 
 
 
 
 
\needspace{5cm} \section{1000 způsobů jak} \label{1000 způsobů jak} \katabox{Činnosti (3-5) 
}{libovolný, obvykle se zapojí všichni 
}{volný, dokud nedojdou nápady} 
 
 
Nebo taky \textbf{Tisíc způsobů jak}. Hráči hledají různé způsoby, jak vykonat činnost, například \odkaz{výměna žárovky}{výměna žárovky}. 
 
 
\subsection{Průběh}  
\odkaz{MC}{mc} vybere od publika činnosti. Poté se hráči postaví do půlkruhu a střídají se. Každý hráč, který vyběhne, předvede volně nějaký svůj alternativní způsob, jak vykonat zadanou činnost. Obvykle MC po chvíli činnost vymění, zadá jinou. 
 
\subsection{ Varianty } \begin{itemize}
\item  \odkaz{Věty, které by neměly zaznít}{věty, které by neměly zaznít}
\item  \odkaz{Metafory}{metafory}
\item  \odkaz{1000 hitů o...}{1000 hitů o...}
\item  Rekvizita tisíckrát jinak - kde se hráči snaží v~kratičkých črtách využít předmět.
\end{itemize}
 
\subsection{ Běžné chyby } \begin{itemize}
\item Sklouzávání k~záchodovému humoru
\item Přehnané použití dynamitu
\end{itemize}
 
 
 
 
 
\needspace{5cm} \section{4, 2, 1} \label{4, 2, 1} \katabox{Téma nebo prostředí 
}{4 
}{1 min} 
 
 
Nebo taky \textbf{čtyři, dva, jedna}. Improvizátoři ztvární příběh, který opakují se zmenšeným počtem herců, ale při zachování počtu postav. 
 
 
\subsection{ Průběh }  
Hrají 4 hráči na zadané téma víceméně volnou improvizaci, kde ale dávají důraz na pohyb a akci. Tato improvizace je krátká, max 2 minuty. Po skončení diváci vyberou 2 hráče, kteří ten stejný příběh 4 postav zahrají sami (střídají postavy, přeskakují z~jedné do druhé). Posléze diváci vyberou jen jednoho, který celý příběh odehraje naposledy a úplně sám (opět hraje všechny 4 postavy). 
 
Může se plést s~\odkaz{Počet slov}{počet slov}, přezdívanou často '1, 3, 5, 7' 
 
\subsection{ Běžné chyby } \begin{itemize}
\item  Postavy příliš mnoho mluví (děj se nedá zapamatovat)
\item  Úvodní scénka v~počtu 4 herců je příliš dlouhá
\end{itemize}
 
\subsection{ Podobné kategorie }  
\begin{itemize}
\item  \odkaz{Polovina času}{polovina času}
\end{itemize}
 
 
 
